\chapter{Federación con capas personalizadas y salidas heterogéneas (FedPerH)} \label{cap3}

\vspace*{5mm}

En esta sección vamos a proponer una solución al problema de clasificación multi-etiqueta en entornos federados. Veremos cómo trataremos los datos del dataset para adaptarlos a nuestro entorno federado, seleccionaremos la herramienta para el sistema federado, el modelo a usar y las modificciones pertinentes. Por último definiremos las métricas elegidas para evaluar la bondad del sistema.

\section{Datos de entrenamiento}

Nuestro objetivo es tener un número de distintos conjuntos de datos con sus propias etiquetas no compartidas con los demás, buscamos simular distintas tareas a entrenar que sean específicas para cada cliente. Para ello no solo vamos a separar los diferentes ejemplos por cada cliente sino que también tenemos que hacer una discriminación por etiquetas.

Como base vamos a usar el conjunto de imágenes multi-etiqueta Corel-5k, que tiene 5000 imágenes de 50 categorías y un diccionario de 260 palabras o clases ~\url{https://www.kaggle.com/datasets/parhamsalar/corel5k}. A partir de él crearemos 4 conjuntos de datos con 10 etiquetas diferentes cada uno ya que vamos a usar 4 clientes en nuestra federación, para hacer la separacion por etiquetas vamos a ordenar las 40 que más se repiten, creamos cuatro grupos de 10 etiquetas de manera que el número de ejemplos de cada grupo sea similiar. Una vez tenemos los 4 grupos seleccionamos del dataset total los ejemplos correspondientes a las etiquetas de cada grupo, por ejemplo si en nuestro grupo tenemos \"agua\" cogemos todos los elementos del conjunto de datos que estén etiquetados con \"agua\", luego eliminamos de esos ejemplos todas las etiquetas que no crrespondan a su grupo. Hacemos esto para los 4 subsets y así conseguimos separar las imágenes de manera que cada conjunto solo contenga etiquetas propias y que el número de ejemplos sea parecido.

Una vez hemos conseguido fabricar unos conjuntos de datos que nos permitan simular una situación real en la que tenemos conjuntos de datos distintos que quieren entrenar un mismo modelo lo siguiente es preparar el sistema que va a procesar esos datos.

\section{Arquitectura de la solución}

Vamos a montar un sistema que nos permita comparar el rendimiento de un clasificador multi-etiqueta en un entorno federado y uno centralizado usando los mismos parámetros para ver si existe mejora en la tarea al pertenecer a un grupo colaborativo de entrenamiento. Tendremos 4 clientes y cada uno realizará una tarea de clasificación distinta ya que cada conjunto de datos va a tener etiquetas distintas. Para que realicen esta tarea vamos a usar modelos personalizados que contarán con una parte comun que es la que se compartirá con el servidor y un conjunto de capas personalizadas que quedarán en los clientes. De esta manera parte del modelo tendrá una función general y compartida entre los demás nodos y cada uno de ellos entrenará la parte específica a su tarea de manera local. Los clientes compartirán la parte común del modelo y el servidor se encargará de realizar la agregación y de devolverle esa parte actualizada todos ellos.

\subsection{FedPer}

Para evitar la implementación de todo el sistema desde cero vamos a usar de base el entorno de trabajo para aprendizaje federado Flower ~\cite{} que nos provee de herramientas de simulación para este tipo de sistemas. Más en concreto haremos uso del baseline FedPer, una implementación de ~\cite{} el cual usa estas capas federadas para una tarea de clasificación multi-clase. 

\subsection{Resnet}

\subsection{Modificaciones}
%Backbone Resnet

%Modificaciones para multi-etiqueta

\section{Métricas}





